\section{Methodology}
\label{sec:method}

We detail the mathematical foundations and algorithmic formulation of the four primary architectural pillars comprising NeuralOps.

\subsection{Pillar 1: Hardware-Aware SRAM-Fused SVD Kernel}
\label{sec:fused_kernel}

Consider an input sequence tensor $X \in \mathbb{R}^{B \times L \times d_{\text{in}}}$ and factorized weight matrices $V \in \mathbb{R}^{r \times d_{\text{in}}}$ and $U \in \mathbb{R}^{d_{\text{out}} \times r}$ representing a rank-$r$ projection ($r \ll \min(d_{\text{in}}, d_{\text{out}})$). Standard implementations decompose the forward pass into two discrete kernel launches:
\begin{align}
Z &= X V^\top \in \mathbb{R}^{B \times L \times r} \quad \xrightarrow{\text{write to HBM}} \text{DRAM}, \\
Z_{\text{norm}} &= \operatorname{RMSNorm}(Z, w) \quad \xrightarrow{\text{read/write}} \text{DRAM}, \\
Y &= Z_{\text{norm}} U^\top + b \in \mathbb{R}^{B \times L \times d_{\text{out}}} \quad \xrightarrow{\text{write to HBM}} \text{DRAM}.
\end{align}

\begin{table}[h]
\centering
\small
\begin{tabularx}{\columnwidth}{lX}
\toprule
\multicolumn{2}{l}{\textbf{Algorithm 1: In-SRAM Fused SVD Execution Tile}} \\
\midrule
\textbf{Input:} & $X \in \mathbb{R}^{M \times d_{\text{in}}}$, $V \in \mathbb{R}^{r \times d_{\text{in}}}$, $U \in \mathbb{R}^{d_{\text{out}} \times r}$, scale $w$, bias $b$. \\
\textbf{Output:} & $Y \in \mathbb{R}^{M \times d_{\text{out}}}$. \\
\midrule
1: & Partition total tokens $M = B \cdot L$ into tiles of size $\mathcal{B}_M$. \\
2: & \textbf{Parallel For} each block $m \in \{0, \dots, \lceil M / \mathcal{B}_M \rceil - 1\}$: \\
3: & \quad Load tile $\mathbf{X}_m \in \mathbb{R}^{\mathcal{B}_M \times d_{\text{in}}}$ into SRAM registers. \\
4: & \quad Compute in-register down-projection: $\mathbf{Z}_m = \mathbf{X}_m V^\top$. \\
5: & \quad Compute in-register variance: $\sigma_m^2 = \frac{1}{r} \sum_{k=1}^r \mathbf{Z}_{m, \cdot, k}^2$. \\
6: & \quad Apply RMSNorm: $\mathbf{Z}_{\text{norm}, m} = \frac{\mathbf{Z}_m}{\sqrt{\sigma_m^2 + \epsilon}} \odot w$. \\
7: & \quad Compute up-projection: $\mathbf{Y}_m = \mathbf{Z}_{\text{norm}, m} U^\top + b$. \\
8: & \quad Write tile $\mathbf{Y}_m$ directly to Global Memory (HBM). \\
\bottomrule
\end{tabularx}
\end{table}

For total tokens $M = B \cdot L$, the standard workflow transfers $M \cdot r \cdot \text{sizeof}(\text{dtype})$ bytes to and from HBM multiple times.

In NeuralOps, we formulate an in-SRAM tiled kernel that retains intermediate low-rank activations $\mathbf{Z} \in \mathbb{R}^{\mathcal{B}_M \times r}$ entirely within on-chip register files.

\textbf{Analytical Autograd Backward Formulation:} During backpropagation, given output gradient $G_Y = \frac{\partial \mathcal{L}}{\partial Y} \in \mathbb{R}^{M \times d_{\text{out}}}$, the gradient with respect to normalized features $Z_{\text{norm}}$ is $G_{Z_{\text{norm}}} = G_Y U \in \mathbb{R}^{M \times r}$.

Let $\hat{z} = \frac{z}{\text{rms}(z)}$ where $\text{rms}(z) = \sqrt{\frac{1}{r} \sum_{j=1}^r z_j^2 + \epsilon}$, and let $g_{\hat{z}} = G_{Z_{\text{norm}}} \odot w$. The exact analytical derivative with respect to unnormalized intermediate state $z$ is:
\begin{equation}
\frac{\partial \mathcal{L}}{\partial z} = \frac{1}{\text{rms}(z)} \left( g_{\hat{z}} - \hat{z} \odot \frac{1}{r} \sum_{j=1}^r (g_{\hat{z}, j} \cdot \hat{z}_j) \right).
\end{equation}

The input and parameter gradients are computed directly as:
\begin{align}
\nabla_X \mathcal{L} &= \left(\frac{\partial \mathcal{L}}{\partial z}\right) V, \quad \nabla_V \mathcal{L} = \left(\frac{\partial \mathcal{L}}{\partial z}\right)^\top X, \\
\nabla_U \mathcal{L} &= G_Y^\top Z_{\text{norm}}, \quad \nabla_b \mathcal{L} = \sum_{m=1}^M G_{Y, m}.
\end{align}

\begin{proposition}[Arithmetic Intensity Bound]
The fused SVD Triton kernel improves arithmetic intensity $\mathcal{I} = \frac{\text{FLOPS}}{\text{Bytes Transferred}}$ by a factor of:
\begin{equation}
\kappa = \frac{2 M r (d_{\text{in}} + d_{\text{out}})}{M (d_{\text{in}} + d_{\text{out}} + 4r) \cdot \text{sizeof}(\text{dtype})} \approx \frac{2r}{\text{sizeof}(\text{dtype})},
\end{equation}
achieving $2.5\times$ acceleration over unfused sequential execution on A100 GPUs.
\end{proposition}


\subsection{Pillar 2: Riemannian Optimization on the Stiefel Manifold}
\label{sec:riemannian}

Let the Stiefel manifold of $p$-frames in $\mathbb{R}^n$ ($n \ge p$) be defined as:
\begin{equation}
\operatorname{St}(p, n) = \left\{ W \in \mathbb{R}^{n \times p} \;\middle|\; W^\top W = I_p \right\}.
\end{equation}

When optimizing low-rank projection matrices $U \in \operatorname{St}(r, d_{\text{out}})$ and $V \in \operatorname{St}(r, d_{\text{in}})$, standard Euclidean gradient updates $W_{t+1} = W_t - \eta G_t$ violate the orthogonality constraint $W^\top W = I$, causing progressive spectral collapse and ill-conditioned forward activations.

\begin{figure}[t]
\centering
% \includegraphics[width=0.48\textwidth]{figure2_stiefel_manifold.png}
\caption{\textbf{Riemannian Geometry on $\operatorname{St}(p, n)$.} The Euclidean ambient gradient $G$ is projected onto tangent space $T_W \operatorname{St}(p, n)$ yielding $\operatorname{grad}_R$, followed by a QR retraction $\mathcal{R}_W(-\eta \operatorname{grad}_R)$ with diagonal sign alignment.}
\label{fig:stiefel_geom}
\end{figure}

\textbf{Riemannian Gradient Projection:} The tangent space at point $W \in \operatorname{St}(p, n)$ is given by:
\begin{equation}
T_W \operatorname{St}(p, n) = \left\{ \xi \in \mathbb{R}^{n \times p} \;\middle|\; W^\top \xi + \xi^\top W = 0 \right\}.
\end{equation}

Given the ambient Euclidean gradient $G = \nabla_W \mathcal{L}$, the canonical Riemannian gradient $\operatorname{grad}_R \in T_W \operatorname{St}(p, n)$ under the Euclidean metric is:
\begin{equation}
\label{eq:stiefel_grad}
\operatorname{grad}_R = \Pi_{T_W \operatorname{St}}(G) = G - W \left(\frac{W^\top G + G^\top W}{2}\right) = G - W \operatorname{sym}(W^\top G),
\end{equation}
where $\operatorname{sym}(A) = \frac{1}{2}(A + A^\top)$. For wide projection matrices ($n < p$), the dual projection holds:
\begin{equation}
\operatorname{grad}_R = G - \operatorname{sym}(G W^\top) W.
\end{equation}

\textbf{QR Retraction with Phase Correction:} To map tangent updates back onto $\operatorname{St}(p, n)$ without expensive matrix exponential calculations, we apply the metric-preserving QR retraction $\mathcal{R}_W: T_W \operatorname{St} \to \operatorname{St}$:
\begin{equation}
Q, R = \operatorname{qr}(W - \eta \operatorname{grad}_R),
\end{equation}
\begin{equation}
W_{t+1} = \mathcal{R}_W(-\eta \operatorname{grad}_R) = Q \operatorname{diag}(\operatorname{sgn}(\operatorname{diag}(R))),
\end{equation}
where $\operatorname{sgn}(x) = 1$ if $x \ge 0$ else $-1$, ensuring positive diagonal elements in $R$ and preventing unphysical phase flips during optimization.

\textbf{Riemannian Momentum Buffering:} We formulate Riemannian momentum accumulation by parallel-transporting velocity vectors via tangent projection:
\begin{equation}
B_{t+1} = \Pi_{T_{W_{t+1}} \operatorname{St}}\left( \beta B_t + (1 - \beta) \operatorname{grad}_R(W_{t+1}) \right).
\end{equation}


\subsection{Pillar 3: Information-Theoretic Dynamic Rank \& Spectral Residual LoRA}
\label{sec:dynamic_rank}

Uniform rank allocation across transformer layers is sub-optimal: lower layers capture low-frequency syntactical features with concentrated singular spectra, whereas upper layers process diffuse contextual representations requiring higher dimensional capacity.

\textbf{Singular Entropy Formulation:} For a weight or empirical activation covariance matrix $C_l \in \mathbb{R}^{d \times d}$ at layer $l \in \{1, \dots, L\}$, let its singular values be $\Sigma_l = (\sigma_1, \sigma_2, \dots, \sigma_d)$ with $\sigma_1 \ge \sigma_2 \ge \dots \ge \sigma_d \ge 0$. We define the normalized spectral distribution:
\begin{equation}
p_i = \frac{\sigma_i}{\sum_{j=1}^d \sigma_j + \epsilon}.
\end{equation}

The \textbf{Singular Entropy} $H(\Sigma_l)$ is defined as:
\begin{equation}
\label{eq:singular_entropy}
H(\Sigma_l) = -\sum_{i=1}^d p_i \ln(p_i + \epsilon) \in [0, \ln d].
\end{equation}

Normalized entropy $\bar{H}_l = \frac{H(\Sigma_l)}{\ln d} \in [0, 1]$ directly quantifies representation richness: $\bar{H}_l \to 0$ corresponds to rank-1 redundancy, while $\bar{H}_l \to 1$ indicates a uniform spectrum.

\textbf{Water-Filling Dynamic Rank Allocation:} Given a total parameter budget $R_{\text{total}} = \sum_{l=1}^L r_l$, we solve the constrained allocation problem:
\begin{equation}
r_l = \operatorname{clip}\left( \left\lfloor R_{\text{total}} \cdot \frac{H(\Sigma_l)}{\sum_{k=1}^L H(\Sigma_k)} \right\rceil, r_{\min}, r_{\max} \right),
\end{equation}
subject to physical dimension clamping $r_l \le \min(d_{\text{in}}^{(l)}, d_{\text{out}}^{(l)})$.

\begin{figure}[t]
\centering
% \includegraphics[width=0.48\textwidth]{figure3_res_svd.png}
\caption{\textbf{Spectral Residual LoRA (Res-SVD) Decomposition.} The dominant singular subspace $U_r, V_r$ is frozen, while adapter matrices $A, B$ are initialized directly from the discarded singular tail $\Sigma_{r+1 : r+k}$.}
\label{fig:res_svd_diag}
\end{figure}

\textbf{Spectral Residual LoRA (Res-SVD):} When compressing a dense matrix $W \in \mathbb{R}^{d_{\text{out}} \times d_{\text{in}}}$ to base rank $r$ with auxiliary PEFT adapter rank $k$, standard LoRA initializes adapters randomly (e.g., $B=0, A \sim \mathcal{N}$), discarding the residual spectrum.

In Res-SVD, we compute the full SVD: $W = \sum_{i=1}^d \sigma_i u_i v_i^\top$. We partition the singular spectrum into base factors ($1 \dots r$) and residual adapter factors ($r+1 \dots r+k$):
\begin{equation}
W_{\text{base}} = \sum_{i=1}^r \sigma_i u_i v_i^\top \quad (\text{frozen, } \operatorname{requires\_grad}=\text{False}),
\end{equation}
\begin{equation}
A = \sqrt{\frac{\Sigma_{\text{tail}}}{\gamma}} V_{\text{tail}}^\top \in \mathbb{R}^{k \times d_{\text{in}}}, \quad B = U_{\text{tail}} \sqrt{\frac{\Sigma_{\text{tail}}}{\gamma}} \in \mathbb{R}^{d_{\text{out}} \times k},
\end{equation}
where $\Sigma_{\text{tail}} = \operatorname{diag}(\sigma_{r+1}, \dots, \sigma_{r+k})$ and $\gamma$ is the LoRA scaling factor $\frac{\alpha}{k}$.

\begin{theorem}[Zero-Shot Spectral Optimality]
Let $W_{\text{approx}} = W_{\text{base}} + \gamma B A$. The Res-SVD initialization satisfies:
\begin{equation}
\|W - W_{\text{approx}}\|_F^2 = \sum_{i=r+k+1}^d \sigma_i^2 < \sum_{i=r+1}^d \sigma_i^2 = \|W - W_{\text{base}}\|_F^2,
\end{equation}
guaranteeing strictly minimal initial reconstruction error compared to standard zero-initialized LoRA.
\end{theorem}


\subsection{Pillar 4: Data-Dependent Decay-Gated Recurrent Linear Attention}
\label{sec:decay_attention}

To prevent associative memory state explosion across long context horizons ($N \gg 4096$), we introduce an input-conditioned exponential decay gate into SVD Linear Attention.

For token representation $x_t \in \mathbb{R}^d$, the scalar or channel-wise decay factor $\lambda_t$ is computed via:
\begin{equation}
\label{eq:decay_gate}
\lambda_t = \exp\left(-\operatorname{softplus}\left(\mathbf{w}_{\text{gate}}^\top x_t + b_{\text{gate}}\right)\right) \in (0, 1].
\end{equation}

Because $\operatorname{softplus}(u) = \ln(1 + e^u) > 0$ for all $u \in \mathbb{R}$, $\lambda_t$ is unconditionally bounded in $(0, 1]$, strictly precluding expansive state divergence ($\lambda_t > 1$).

\textbf{Recurrent State Update Equations:}
Given kernel feature maps $\phi(Q_t) = \operatorname{ELU}(Q_t) + 1 \in \mathbb{R}^{H \times 1 \times d_k}$ and $\phi(K_t) = \operatorname{ELU}(K_t) + 1 \in \mathbb{R}^{H \times 1 \times d_k}$:
\begin{align}
S_t &= \lambda_t S_{t-1} + \phi(K_t)^\top V_t \in \mathbb{R}^{H \times d_k \times d_v}, \label{eq:state_S} \\
Z_t &= \lambda_t Z_{t-1} + \phi(K_t)^\top \mathbf{1} \in \mathbb{R}^{H \times d_k \times 1}, \label{eq:state_Z} \\
Y_t &= \frac{\phi(Q_t) S_t}{\phi(Q_t) Z_t + \epsilon} \in \mathbb{R}^{H \times 1 \times d_v}. \label{eq:output_Y}
\end{align}

\begin{theorem}[Steady-State Associative Memory Boundedness]
Assuming bounded inputs $\|\phi(K_t)\|_F \le M_K$ and $\|V_t\|_F \le M_V$, and a uniform minimum decay rate $\lambda_t \le \lambda_{\max} < 1$, the associative memory Frobenius norm is unconditionally bounded for arbitrary sequence length $N \to \infty$:
\begin{equation}
\|S_N\|_F \le \sum_{j=1}^N \lambda_{\max}^{N-j} M_K M_V \le \frac{M_K M_V}{1 - \lambda_{\max}} < \infty.
\end{equation}
\end{theorem}
